\section{Discussion}

On the first question, a general-purpose LLM that received raw 12-lead recordings as files and processed them with code it wrote itself
classified PTB-XL ECGs into the five diagnostic superclasses with a macro AUROC of \MainAUROC{} on the benchmark test
fold, without ECG-specific training, retrieval, labelled examples or class prevalences in the prompt, and every record
received an answer. It exceeded the control that used only age and sex by $\CodeVsAgeSex$ and is
\GapToBestPublishedAbs{} below the best published result with the same split, label
definition and metric, reached by a supervised ensemble~\citep{Strodthoff2021} and by MultiresNet~\citep{shi2023multires},
both trained on the labelled training folds. In an
informal comparison, the result exceeds the best zero-shot ECG--language model~\citep{cai2025supreme} by
$\GapToBestZeroShot$; that model was pretrained on hundreds of thousands of ECGs with report-derived labels. In other
studies, general-purpose LLMs that were given ECG images or numeric series as text, including recent models, have scored near
chance or missed most critical findings~\citep{liu2024teach,seki2025assessing,ahammad2026ecgimage,soubh2026performance,gao2026luminaecg,liu2025picture,langer2025opentslm}.

On the second question, the transcripts, analysed post hoc, show how the model used the sandbox. It wrote code to filter the signal and detect
beats in every record and to form median beats in almost every record, but it rarely printed the full measurement
table that the prompt asked for. Its code usually printed a heart rate, but for most of the PR intervals, QT or QTc
intervals and axes that the model stated, no value was printed under the measurement's own label, and a Sokolow--Lyon sum
was almost never printed. The printed median-beat samples suggest that the model read these
intervals and axes from samples of the median beats it had built. The sandbox thus appears to have served less to
replace reading with computation than to give the model median beats of its own making from the full recording, in the
leads, segments and resolution it chose. Our data cannot show why the model departed from the prompt, but two
observations bear on it. The values printed most often under their own label were those that follow directly from beat
detection (heart rate in \PHPrintedHrPct\% of records, QRS duration in \PHPrintedQrsPct\%). The PR and QT intervals,
printed in \PHPrintedPrPct\% and \PHPrintedQtcPct\%, require locating the onset of the P wave and the end of the T wave,
which is harder to do automatically~\citep{martinez2004wavelet}; in a challenge on QT measurement, the best automated
method erred more than twice as much as the best manual one~\citep{moody2006qt}. The axis and the Sokolow--Lyon voltage
are simple sums of amplitudes that the model could form from a printed per-lead amplitude table (\PHLeadAmplitudePct\%
of records). Moreover, the prompt asked for a table of more than 20 measurements, eight of them for every lead, and
LLMs satisfy fewer of an instruction's requirements as their number grows~\citep{jiang2024followbench}. Because the
model sees only what its code prints, a value computed but not printed could not inform the answer; each stated value
came from printed output or from the model's own reading of printed numbers.

On the third question, the superclass probabilities, analysed post hoc by diagnostic statement, best separated records
with complete
bundle branch blocks, left anterior fascicular block, anterolateral or non-specific ischaemic change or anteroseptal
infarction from records without the superclass (AUROC \PHStmtASMI--\PHStmtCRBBB). They did so least well for
incomplete right bundle branch block and non-specific intraventricular conduction disturbance (\PHStmtIRBBB{} and
\PHStmtIVCD), two statements defined by borderline QRS durations. This measures how well the superclass probability
separated these records, not whether the statement itself was identified. For these two statements, the QRS durations
that the model stated rarely met the prompt's criteria (110--119\,ms in \PHMeasIRBBBCodeCrit\% and 120\,ms or more in
\PHMeasIVCDCodeCrit\% of the records).

The remaining gap to supervised networks is unlikely to reflect missing information in the recordings, since networks
trained on the same recordings reach \BestPublishedAUROC~\citep{Strodthoff2021}. A zero-shot model can apply only the
criteria it is given or knows, and PTB-XL's labels reflect the annotators' practice as well as textbook thresholds. In
post hoc analyses (Supplementary Table~\ref{tab:subgroups}), MI discrimination fell from \PHMILikHundred{} to \PHMILikLow{} as the highest likelihood that the
annotator assigned to a record's MI statements fell from 100 to 15--35. Part of this fall is expected for any
classifier (the age-and-sex control also fell, from \PHMILikHundredAgeSex{} to \PHMILikLowAgeSex), and without
per-record predictions of supervised networks we did not test whether theirs falls less. Among records labelled left
ventricular hypertrophy (LVH) whose key findings stated a Sokolow--Lyon voltage, about half (\PHMeasLVHCodeCrit\%)
stated 3.5\,mV or more (median \PHMeasLVHCodeMedian\,mV); LVH labels may rest on other criteria or on the annotators'
judgement. Supervised models may learn such labelling conventions from the labels, but this is a hypothesis that our data
cannot test. The stated measurements may also reflect the 0.5--40\,Hz filter that the model chose in every record: an
upper cutoff of 40\,Hz, far below the 150\,Hz recommended for diagnostic ECGs in
adults~\citep{kligfield2007recommendations}, can reduce QRS amplitudes and blur QRS onset and offset. In the
development runs, labelled examples attached as files did not measurably improve on the frozen prompt: on the same
100 records, macro AUROC was \DevBoth{} with them and \DevVTwoHundred{} without, a paired difference of
$\DevBothDiff$ ($\DevBothDiffLow$ to $\DevBothDiffHigh$) that was computed after the freeze from the same runs. Nor did
rules for PTB-XL's labelling conventions (Table~\ref{tab:development}), although the samples were small.

Each answer came with the code the model ran, what the code printed and the key findings, and the complete transcript of
every request is part of the released data. These records make the model's work inspectable in a way that class
scores are not: a reader can check how the signal was filtered, how the beats were detected and what was printed, as in
Fig.~\ref{fig:example}. They can also be checked by execution: replayed offline, the model's code reproduced the
printed output of the API's sandbox in all \RPRuns{} code runs, up to floating-point rounding, the order of tied
values and a randomly named directory. This confirms the computations behind each answer, not that the model would write the same code again. The
records document how the median beats were formed more fully than how each stated value was
obtained, because most stated PR intervals, QT or QTc intervals and axes appear to have been read from printed samples
rather than computed, and the stated values were rounded and sometimes differed from printed ones. On the second part of the third question, the stated values, compared post hoc,
were consistent with the cardiologists' statements for rate, rhythm and QRS duration in
complete bundle branch block, partly for the PR interval in first-degree atrioventricular block and the axis in left
anterior fascicular block, for about half of the LVH records that stated a Sokolow--Lyon voltage, and poorly for
incomplete right bundle branch block and non-specific intraventricular conduction disturbance. The consistency for rate,
rhythm and complete bundle branch block suggests that the key findings were specific to each recording. However,
consistency with categorical statements is not accuracy: the stated values were not compared with reference
measurements, such as those that ECG software provides for PTB-XL~\citep{Strodthoff2023ptbxlplus}. The reasoning
returned by the API is a summary, and we did not test whether the stated measurements determined the probabilities.

For practice, the approach needs no labelled data for model training, although the prompt was developed on
\NDevRecordsUsed{} labelled validation records; it needs only a general-purpose LLM with a Python sandbox and a prompt.
In the configuration tested here, the approach is currently not ready for clinical use. Discrimination was below that of trained networks, the superclass task is coarse and
leaves out rhythm diagnoses, the probabilities would need recalibration, and the model did not follow its measurement
instructions reliably. Inference currently costs about USD \CostPerRecord{} per ECG at batch prices (USD \CostPerThousand{}
per thousand; Supplementary Section~\ref{app:cost}), far more than running a trained network, and the batch interface returns results after a delay. Sending raw
ECGs to an external service requires data-governance arrangements, and clinical use would require prospective
validation and review of errors by cardiologists. A potential advantage, not tested here, is that the same model can be
redirected to other criteria or tasks by changing the prompt rather than by retraining.

Beyond these limitations of the approach, the study itself has limitations. It used one public dataset. PTB-XL and its labels have been public since
2020~\citep{Wagner2020} and derived measurements since 2023~\citep{Strodthoff2023ptbxlplus}, all before the model's
training-data cutoff (June 2026~\citep{anthropic2026models}). The model knew that it was classifying PTB-XL records,
because the prompt names PTB-XL's superclasses, and it had some knowledge of the dataset: the API returned a non-empty
reasoning summary for \PHSummariesNonEmptyN{} of \NTest{} records, and \PHThinkPTBXLN{} of these
(\PHThinkPTBXLNonEmptyPct\%) mentioned PTB-XL by name (post hoc), mostly its labelling conventions. In the example shown in Fig.~\ref{fig:example}, the model reasoned
that voltage-only LVH is often labelled as LVH in datasets like PTB-XL (reasoning summary in Supplementary
Section~\ref{app:transcript}). The
files carried anonymous names and no identifiers, so the model could not look a record up by name, but we cannot exclude
that exposure to PTB-XL or to work that used it during training affected the result. The setting is zero-shot in a
narrow sense: the model's parameters were not updated and the prompt contained no labelled examples or class
prevalences, but the prompt defines the superclasses as PTB-XL does, and the configuration was chosen on labelled
validation records. In this configuration, the test fold was read once, by one model from one vendor. No configuration of this model without the sandbox is reported
here, so the study does not show how much of the result is due to the sandbox itself. Outputs are not deterministic, no repeat runs were made, and the confidence
intervals reflect the sampling of records only. Variants of the frozen prompt were compared on 100 or 200 records, where
the 95\% CIs of the paired differences had half-widths of about \DevFreshVFiveDiffHalfWidth{} (200 records) and
\DevBothDiffHalfWidth--\DevStatementsDiffHalfWidth{} (100 records), so differences of about that size or smaller could
not be reliably detected. The diagnostic criteria in the prompt are condensed and in places simplified (Supplementary
Table~\ref{tab:criteria}): poor R-wave progression is listed for MI although it is not a diagnostic criterion, long
QTc is listed under STTC, as PTB-XL classifies it, although it is not specific to ischaemia, the threshold for
non-specific intraventricular conduction disturbance is higher than the recommended one, and the age- and
sex-specific thresholds for ST elevation in V2--V3 and most criteria for RVH are not given. Whether criteria stated as
in the guidelines would improve discrimination was not tested; it would need new development runs. It should be
noted that the comparisons with published results (Fig.~\ref{fig:benchmark}a)
are informal, because inputs (sampling rates; age and sex,
which only this model received) and in some cases splits differ and per-record predictions were not available. The
measurements were extracted from free text with regular expressions, and the statements they were compared with are
categorical labels of a single annotation process that, following the benchmark definition, include statements of any
likelihood. The analysis of what the model did in the sandbox, the per-statement results, the subgroups, the calibration
analysis, the MI likelihood strata, the patient-level bootstrap, the age-and-sex control on the validation fold, the
measurement analysis, the count of PTB-XL mentions in the reasoning summaries and the local replay of the model's code
were added after the test run.

Evaluation on ECGs recorded after the model's training-data cutoff would address both generalization and possible exposure
to PTB-XL. Comparing the stated measurements with reference measurements, diagnostic criteria stated as in the guidelines,
prompts or tools that make the model compute
its measurements in code, having cardiologists review the transcripts of errors, stratifying the predictions of a
supervised network by annotation likelihood, extending the approach to PTB-XL's finer diagnostic and rhythm
tasks~\citep{Strodthoff2021}, repeating runs and testing LLMs from other vendors would show how far the approach
extends and how far its written work can be trusted.

In conclusion, given a Python sandbox and raw 12-lead recordings, a general-purpose LLM processed each ECG with its own
code and classified the PTB-XL superclasses with a macro AUROC of \MainAUROC{} without training on labelled ECGs,
between published zero-shot ECG--language models and supervised networks. It used code mainly to prepare median
beats, and its transcripts suggest that it read most of the intervals and axes it stated from samples it printed;
every answer came with the code the model ran, what the code printed and the measurements the model stated.
